\section{Discussion}
\label{sec:discussion}

\paragraph{No detectable directional contribution.}
A wind-weighted fire-exposure feature looks like a direct encoding of smoke transport. In our
data, once seed noise is averaged out, it is at most a few thousandths of average precision better
than the same exposure combined with a wind from another time, with directions permuted within
the month, or with no direction at all but the same missingness, and not consistently so. Within
the precision of an eleven-season record, the alignment between the reported wind and the fire
positions adds little or nothing that these controls do not also provide.
This is a statement about this feature, this target and this evaluation, not a proof that
direction never matters for smoke.

\paragraph{Single fits of tree models are not enough.}
Our first, pre-specified analysis compared single LightGBM fits and found a large, significant
directional gain, which several placebos then appeared to reproduce. We first read that as a
mechanism: fire exposure behaving differently in calm, stagnant air. It was neither. The
differences between variants were mostly differences between random seeds, concentrated in one
extreme test year that the trees had to extrapolate to. A paired bootstrap does not catch this,
because it resamples test rows for fixed fitted models and says nothing about how much the fitted
models themselves would change with another seed. For rare, clustered events evaluated across
seasons, we recommend at least a seed ensemble, or a report of the spread over seeds, before
comparing feature sets.

\paragraph{Placebo design.}
The obvious placebo, rotating the wind, is not a clean null: it keeps the current wind axis and
the exact missingness pattern of the directional feature. Time-shifted winds are better but not
independent of the current wind. The 168\,h-shifted direction still has a mean cosine similarity
of \pv{shiftdiag.168h.cos} with the current one, and even 720\,h earlier it is
\pv{shiftdiag.720h.cos}, because monsoon flow persists for weeks; its missingness matches the
current report in only \pv{shiftdiag.168h.mask}\% of hours. Permuting directions within the month
keeps the missingness exactly and lowers the similarity to \pv{shiftdiag.perm1.cos}. We recommend
mask-preserving placebos of this kind, together with a missingness-matched non-directional
control, for any study that credits wind-aware fire features.

\paragraph{Why direction may not matter here.}
Several explanations are plausible, and our data cannot separate them. The airport's surface wind
is a poor proxy for transport in the boundary layer, especially at night and near the coast,
where sea breezes dominate several of our stations. Smoke from large peat fires is regional: in
strong seasons the whole region is covered, so the direction of the nearest fires matters less
than the total. The 300\,km and 72\,h window is long relative to how fast surface winds change.
The persistent-flow variant was meant to address the last two points; it was run only as single
fits and is not interpretable on its own. Trajectory- or reanalysis-based transport features are
the natural next test.

\paragraph{Model families depend on the split.}
Logistic regression on station state beats every gradient-boosted model under leave-one-year-out,
yet on the operational split gradient boosting is ahead (AP \pv{eop24.AP.D} for counts against
\pv{eop24.AP.lgC}). The ranking depends on whether the test period lies inside or outside the
range seen in training. Leave-one-year-out trains on later years to score earlier ones, so it
measures generalisation across seasons, not operational skill; for deployment decisions the
forward-in-time split is the more relevant one.

\paragraph{Implications for Pavana.}
The deployed model uses detection counts plus the directional exposure. Our results support
keeping the counts and suggest the directional term can be dropped with no measurable loss. We
report this as found and have not changed the deployed system as part of this study.

\paragraph{Limitations.}
Visibility with a haze or smoke code is an observer-reported marker, not a smoke measurement; HZ
is also used for dust and urban pollution, coding differs between countries and changes over
time, and IEM applies little quality control. Several Indonesian stations have long reporting
gaps before 2021, and Palangka Raya, in the most fire-affected province, is absent. We use one
satellite sensor and processed archive detections rather than the near-real-time records that
were available operationally; our latency shift moves detection times but cannot reproduce
differences between the two versions. The exposure kernel is fixed from the deployed model; other
distance or age scales might behave differently. Detections are not confirmed fires, and their
absence under thick smoke or cloud biases exposure downward exactly when haze is worst. Bootstrap
intervals treat the fitted models as fixed, and five seeds leave residual seed noise; intervals
with all-station blocks are wider than the station-level ones, and eleven years are far from
eleven independent episodes. Finally, most controls in this paper were added after the first
results (Appendix~\ref{app:protocol}); they were added to test the first result, and they removed
it.

\section{Conclusion}

Across eleven years and twelve airports in Borneo and Sumatra, satellite fire counts modestly
improve 24\,h forecasts of haze-limited visibility, but wind-directional fire exposure adds at
most a few thousandths of average precision, inconsistently, over a time-shifted or permuted-wind
placebo or a non-directional exposure with the same missingness. A single-fit comparison would have suggested a large
directional gain; it was seed noise in one extreme season. Studies of wind-aware fire features
should average over seeds and compare against mask-preserving placebos before attributing skill
to direction.
